\ELhold

\ELsection{Course Outline}{Course Outline}

\begin{ELunit}

\ELfigure{m01-course-outline}{}

\end{ELunit}

\ELhold

\ELsection{Vector Analysis}{Vector Analysis}

\begin{ELunit}

\begin{itemize}
\tightlist
\item
  \textbf{Vector analysis} is a mathematical tool with which the expression and understanding of electromagnetic concepts become simple.
\item
  Quantities in electromagnetics are divided into the following two groups.

  \begin{itemize}
  \tightlist
  \item
    \textbf{Scalar:} a quantity that is completely specified by its magnitude (positive or negative). For example: charge, electric current, energy, \ldots{}
  \item
    \textbf{Vector:} a quantity that is described by a magnitude and a direction. For example the electric and magnetic field intensity, \ldots{}
  \end{itemize}
\item
  Both of the above quantities can be functions of position and time.
\end{itemize}

\end{ELunit}

\ELnextnum{1-1}

\begin{ELexample}{}

\begin{itemize}
\tightlist
\item
  In your opinion, of what type are the following quantities?

  \begin{itemize}
  \tightlist
  \item
    Temperature
  \item
    Power
  \item
    Force
  \item
    Potential difference
  \item
    Velocity
  \end{itemize}
\item
  Apart from the above, which other scalar and vector quantities in physics can you name?
\end{itemize}

\end{ELexample}

\begin{ELunit}

\ELfigure{m01-vector-analysis-map}{}

\begin{itemize}
\tightlist
\item
  How can two vector quantities be added or subtracted graphically?
\item
  How can vectors be multiplied? What is the meaning of this product? Where is it used?
\end{itemize}

\ELdisp{\vect{A}\ELbr =2\uvec{x}+3\uvec{y}+\uvec{z}\qquad\ELbr  \vect{B}\ELbr =\uvec{x}-2\uvec{y}+3\uvec{z}\qquad\ELbr  \vect{C}\ELbr =-\uvec{x}+4\uvec{y}-\uvec{z}}

\begin{itemize}
\tightlist
\item
  Find the magnitude of the projection of vector \ELim{\vect{A}} on vector \ELim{\vect{B}}.
\item
  What is the angle between vectors \ELim{\vect{A}} and \ELim{\vect{B}}?
\item
  What is the area of the triangle formed by the two vectors \ELim{\vect{A}} and \ELim{\vect{B}}?
\item
  Find the unit vector normal to the plane containing the two vectors \ELim{\vect{A}} and \ELim{\vect{B}}.
\item
  What is the volume of the parallelepiped formed by the 3 vectors \ELim{\vect{A}}, \ELim{\vect{B}} and \ELim{\vect{C}}?
\end{itemize}

\end{ELunit}

\ELhold

\ELsection{Vector Algebra}{Vector Algebra}

\begin{ELunit}

\begin{itemize}
\tightlist
\item
  Every vector has a magnitude and a direction.
\end{itemize}

\ELdisp{\vect{A}\ELbr =\uvec{A}A,\qquad\ELbr  \uvec{A}\ELbr =\frac{\vect{A}}{\abs{\vect{A}}}\ELbr =\frac{\vect{A}}{A}\quad\ELbr (\text{unit vector})}

\ELfigure{m01-vector-magnitude}{}

\end{ELunit}

\begin{ELremark}{}

For every vector quantity both the magnitude and the direction must be known. Leaving the direction of a vector quantity unspecified means ignoring part of the information of that quantity.

\end{ELremark}

\ELhold

\ELsection{Addition of Vectors}{Addition of Vectors}

\begin{ELunit}

\begin{itemize}
\tightlist
\item
  Adding two vectors graphically

  \begin{itemize}
  \tightlist
  \item
    The \textbf{parallelogram} rule
  \item
    The \textbf{head-to-tail} rule
  \end{itemize}
\end{itemize}

\ELfigure{m01-vector-sum}{}

\begin{itemize}
\tightlist
\item
  Laws governing the addition of vectors

  \begin{itemize}
  \tightlist
  \item
    \textbf{Commutative law}
    \ELdisp{\vect{A}+\vect{B}\ELbr =\vect{B}+\vect{A}}
  \item
    \textbf{Associative law}
    \ELdisp{\vect{A}+(\vect{B}+\vect{C})\ELbr =(\vect{A}+\vect{B})+\vect{C}}
  \end{itemize}
\end{itemize}

\end{ELunit}

\ELhold

\ELsection{Subtraction of Vectors}{Subtraction of Vectors}

\begin{ELunit}

\begin{itemize}
\tightlist
\item
  Vector subtraction
\end{itemize}

\ELdisp{\vect{A}-\vect{B}\ELbr =\vect{A}+(-\vect{B})}

\ELfigure{m01-vector-difference}{}

\end{ELunit}

\ELhold

\ELsection{Multiplication of Vectors}{Multiplication of Vectors}

\begin{ELunit}

\ELfigure{m01-products-map}{}

\begin{itemize}
\tightlist
\item
  Multiplication of a vector by a scalar:
\end{itemize}

\ELdisp{k\vect{A}\ELbr =\uvec{A}(kA)}

\end{ELunit}

\ELhold

\ELsection{Dot or Scalar Product}{Dot or Scalar Product}

\begin{ELdefinition}{Mathematical definition}

\ELdisp{\vect{A}\cdot\vect{B}\triangleq AB\cos\theta_{AB}}
- It equals the product of the magnitude of one vector and the projection of the other vector on the first.

\end{ELdefinition}

\begin{ELunit}

\ELfigure{m01-dot-product}{\ELim{\theta_{AB}}: the smaller angle between the two vectors when their
tails are joined}

\begin{itemize}
\tightlist
\item
  Properties of the dot product

  \begin{itemize}
  \tightlist
  \item
    The dot product of two vectors is a scalar quantity.
  \item
    It is less than or equal to the product of their magnitudes.
  \item
    Depending on the angle between them it can be positive or negative.
  \item
    If the vectors are perpendicular to each other, it is zero.
  \item
    It is commutative and distributive.
    \ELdisp{\vect{A}\cdot\vect{B}\ELbr =\vect{B}\cdot\vect{A},\qquad\ELbr  \vect{A}\cdot(\vect{B}+\vect{C})\ELbr =\vect{A}\cdot\vect{B}+\vect{A}\cdot\vect{C}}
  \item
    The dot product of a vector with itself
    \ELdisp{\vect{A}\cdot\vect{A}\ELbr =A^2,\qquad\ELbr  A\ELbr =+\sqrt{\vect{A}\cdot\vect{A}}}
  \end{itemize}
\end{itemize}

\end{ELunit}

\ELnextnum{1-2}

\begin{ELexample}{}

Using the concept of the dot product, find the projection of vector \ELim{\vect{A}} on \ELim{\vect{B}}.

\ELfigure{m01-ex2-statement}{}

\begin{ELsolution}

\ELfigure{m01-ex2-projection}{}

Magnitude of the projection of vector \ELim{\vect{A}} on \ELim{\vect{B}}:
\ELdisp{A\cos\theta_{AB}\ELbr =\frac{\vect{A}\cdot\vect{B}}{B}}
Projection of vector \ELim{\vect{A}} on \ELim{\vect{B}}:
\ELdisp{\frac{\vect{A}\cdot\vect{B}}{B}\uvec{B}\ELbr =\left(\frac{\vect{A}\cdot\vect{B}}{B}\right)\frac{\vect{B}}{B}\ELbr =\frac{\vect{A}\cdot\vect{B}}{B^2}\vect{B}}

\end{ELsolution}

\end{ELexample}

\ELnextnum{1-3}

\begin{ELexample}{}

Prove the law of cosines for a triangle using the concept of the dot product.
\ELdisp{C\ELbr =\sqrt{A^2+B^2-2AB\cos\alpha}}

\ELfigure{m01-ex3-triangle}{}

\begin{ELsolution}

\ELfigure{m01-ex3-solution}{}

\ELdisp{\vect{C}\ELbr =\vect{A}+\vect{B}}
\ELdisp{\begin{aligned}C^2=\vect{C}\cdot\vect{C}&=(\vect{A}+\vect{B})\cdot(\vect{A}+\vect{B})\\&=\vect{A}\cdot\vect{A}+\vect{B}\cdot\vect{B}+2\vect{A}\cdot\vect{B}\\&=A^2+B^2+2AB\cos\theta_{AB}\end{aligned}}
\ELdisp{\cos\theta_{AB}\ELbr =\cos(180^\circ-\alpha)\ELbr =-\cos\alpha}
\ELdisp{C^2\ELbr =A^2+B^2-2AB\cos\alpha}

\end{ELsolution}

\end{ELexample}

\ELhold

\ELsection{Cross or Vector Product}{Cross or Vector Product}

\begin{ELdefinition}{}

\ELdisp{\vect{A}\times\vect{B}\triangleq\uvec{n}\abs{AB\sin\theta_{AB}}}

\end{ELdefinition}

\begin{ELunit}

\begin{itemize}
\tightlist
\item
  The magnitude of this product equals the area of the parallelogram formed by the vectors \ELim{\vect{A}} and \ELim{\vect{B}}, and its direction is perpendicular to both vectors.
\end{itemize}

\ELfigure{m01-cross-product}{}

\end{ELunit}

\begin{ELjoin}

\begin{itemize}
\tightlist
\item
  Properties of the cross product

  \begin{itemize}
  \tightlist
  \item
    The cross product of two vectors is a vector quantity.
  \item
    It is not commutative
    \ELdisp{\vect{B}\times\vect{A}\ELbr =-\vect{A}\times\vect{B}}
  \item
    It is distributive
    \ELdisp{\vect{A}\times(\vect{B}+\vect{C})\ELbr =\vect{A}\times\vect{B}+\vect{A}\times\vect{C}}
  \item
    It is not associative
    \ELdisp{\vect{A}\times(\vect{B}\times\vect{C})\neq(\vect{A}\times\vect{B})\times\vect{C}}
  \end{itemize}
\end{itemize}

\ELnextnum{1-4}

\begin{ELexample}{}

Prove the law of sines for a triangle using the concept of the cross product.
\ELdisp{\frac{\sin A}{a}\ELbr =\frac{\sin B}{b}\ELbr =\frac{\sin C}{c}}

\ELfigure{m01-ex4-triangle}{}

\begin{ELsolution}

Area of the triangle:
\ELdisp{\abs*{\tfrac12\vect{a}\times\vect{b}}\ELbr =\abs*{\tfrac12\vect{b}\times\vect{c}}\ELbr =\abs*{\tfrac12\vect{c}\times\vect{a}}}
\ELdisp{ab\sin C\ELbr =bc\sin A\ELbr =ca\sin B}
\ELdisp{\frac{\sin A}{a}\ELbr =\frac{\sin B}{b}\ELbr =\frac{\sin C}{c}}

\end{ELsolution}

\end{ELexample}

\end{ELjoin}

\ELhold

\ELsection{Product of Three Vectors}{Product of Three Vectors}

\begin{ELunit}

\begin{itemize}
\tightlist
\item
  \textbf{Scalar triple product}
\end{itemize}

\ELdisp{\vect{A}\cdot(\vect{B}\times\vect{C})\ELbr =\vect{B}\cdot(\vect{C}\times\vect{A})\ELbr =\vect{C}\cdot(\vect{A}\times\vect{B})}

\begin{itemize}
\tightlist
\item
  Its magnitude equals the volume of the parallelepiped formed by the three vectors \ELim{\vect{A}}, \ELim{\vect{B}} and \ELim{\vect{C}}.

  \begin{itemize}
  \tightlist
  \item
    Area of the base: \ELim{\abs{\vect{B}\times\vect{C}}=\abs{BC\sin\theta_1}}
  \item
    Height: \ELim{\abs{A\cos\theta_2}}
  \item
    Volume of the parallelepiped: \ELim{\abs{ABC\sin\theta_1\cos\theta_2}}
  \end{itemize}
\end{itemize}

\ELfigure{m01-triple-product}{}

\begin{itemize}
\tightlist
\item
  \textbf{Vector triple product}

  \begin{itemize}
  \tightlist
  \item
    Known as the back-cab rule
  \end{itemize}
\end{itemize}

\ELdisp{\vect{A}\times(\vect{B}\times\vect{C})\ELbr =\vect{B}(\vect{A}\cdot\vect{C})-\vect{C}(\vect{A}\cdot\vect{B})}

\end{ELunit}

\ELnextnum{1-5}

\begin{ELexample}{}

The expression \ELim{(\vect{A}\times\vect{B})\times\vect{C}} is equal to which of the following expressions?

\begin{itemize}
\tightlist
\item
  \ELim{\vect{B}(\vect{A}\cdot\vect{C})-\vect{C}(\vect{A}\cdot\vect{B})}
\item
  \ELim{-\vect{B}(\vect{A}\cdot\vect{C})+\vect{C}(\vect{A}\cdot\vect{B})}
\item
  \ELim{-\vect{A}(\vect{C}\cdot\vect{B})+\vect{B}(\vect{C}\cdot\vect{A})}
\item
  \ELim{\vect{A}(\vect{C}\cdot\vect{B})-\vect{B}(\vect{C}\cdot\vect{A})}
\end{itemize}

\begin{ELsolution}

\ELdisp{\begin{aligned}(\vect{A}\times\vect{B})\times\vect{C}&=-\vect{C}\times(\vect{A}\times\vect{B})\\&=-\vect{A}(\vect{C}\cdot\vect{B})+\vect{B}(\vect{C}\cdot\vect{A})\end{aligned}}

\end{ELsolution}

\end{ELexample}

\ELhold

\ELsection{Orthogonal Coordinate Systems}{Orthogonal Coordinate Systems}

\begin{ELunit}

\ELfigure{m01-vector-analysis-map-coord}{}

\begin{itemize}
\item
  What are the principal surfaces in the different coordinate systems?
\item
  How are the unit vectors defined in the different coordinate systems?
\item
  How are the position vector of a point and the distance vector between two points computed in the different coordinate systems?
\item
  Determine the differential length, surface and volume in the different coordinate systems, for use in integration.
\item
  Find the representation of the vector \ELim{\vect{A}=\uvec{r}(3\cos\phi)-\uvec{\phi}2r+\uvec{z}5} in Cartesian coordinates.
\item
  The position of point \ELim{P} in spherical coordinates is \ELim{(8,120^\circ,330^\circ)}. Express the position of this point in cylindrical and Cartesian coordinates.
\item
  To describe the spatial variations of physical quantities, we must be able to describe all points of space uniquely and appropriately.

  \begin{itemize}
  \tightlist
  \item
    This requires the use of an appropriate coordinate system.
  \end{itemize}
\item
  The laws of electromagnetics do not depend on the coordinate system.
\item
  For simplicity of the calculations, a coordinate system that suits the geometry of the problem is used.
\item
  In three-dimensional space a point can be located at the intersection of three surfaces.

  \begin{itemize}
  \tightlist
  \item
    If these three surfaces are mutually perpendicular, we have an orthogonal coordinate system.
  \end{itemize}
\item
  The three orthogonal systems used in this course

  \begin{itemize}
  \tightlist
  \item
    Cartesian (rectangular) coordinates
  \item
    Cylindrical coordinates
  \item
    Spherical coordinates
  \end{itemize}
\end{itemize}

\end{ELunit}

\ELhold

\ELsection{Cartesian Coordinate System}{Cartesian Coordinate System}

\begin{ELunit}

\ELfigure{m01-cartesian-planes}{}

\begin{itemize}
\item
  Every point \ELim{P(x_1,y_1,z_1)} is the intersection of the three planes \ELim{x=x_1}, \ELim{y=y_1} and \ELim{z=z_1}.
\item
  Unit vector \ELim{\uvec{x}} \ELim{\ELto} the unit vector normal to the \ELim{x}-constant plane, in the positive direction
\item
  Unit vector \ELim{\uvec{y}} \ELim{\ELto} the unit vector normal to the \ELim{y}-constant plane, in the positive direction
\item
  Unit vector \ELim{\uvec{z}} \ELim{\ELto} the unit vector normal to the \ELim{z}-constant plane, in the positive direction
\item
  Cross and dot products of the unit vectors
\end{itemize}

\ELdisp{\uvec{x}\times\uvec{y}\ELbr =\uvec{z},\qquad\ELbr  \uvec{y}\times\uvec{z}\ELbr =\uvec{x},\qquad\ELbr  \uvec{z}\times\uvec{x}\ELbr =\uvec{y}}

\ELfigure{m01-cartesian-cycle}{}

\ELdisp{\uvec{x}\cdot\uvec{y}\ELbr =\uvec{y}\cdot\uvec{z}\ELbr =\uvec{z}\cdot\uvec{x}\ELbr =0}
\ELdisp{\uvec{x}\cdot\uvec{x}\ELbr =\uvec{y}\cdot\uvec{y}\ELbr =\uvec{z}\cdot\uvec{z}\ELbr =1}

\end{ELunit}

\begin{ELjoin}

\begin{itemize}
\tightlist
\item
  Representation of an arbitrary vector \ELim{\vect{A}} in Cartesian coordinates
\end{itemize}

\begin{ELimportant}{}

\ELdisp{\vect{A}\ELbr =\uvec{x}A_x+\uvec{y}A_y+\uvec{z}A_z}
\ELdisp{A\ELbr =\sqrt{A_x^2+A_y^2+A_z^2}}

\end{ELimportant}

\end{ELjoin}

\begin{ELunit}

\begin{itemize}
\tightlist
\item
  \textbf{Position vector}

  \begin{itemize}
  \tightlist
  \item
    The position vector of point \ELim{P} is the directed distance from the origin to \ELim{P}.
  \end{itemize}
\end{itemize}

\ELdisp{\vect{R}_1\ELbr =x_1\uvec{x}+y_1\uvec{y}+z_1\uvec{z},\qquad\ELbr  \vect{R}_2\ELbr =x_2\uvec{x}+y_2\uvec{y}+z_2\uvec{z}}

\begin{itemize}
\tightlist
\item
  \textbf{Distance vector}

  \begin{itemize}
  \tightlist
  \item
    The displacement vector from one point to another.
  \end{itemize}
\end{itemize}

\ELdisp{\vect{R}_{12}\ELbr =\vect{R}_2-\vect{R}_1\ELbr =(x_2-x_1)\uvec{x}+(y_2-y_1)\uvec{y}+(z_2-z_1)\uvec{z}}

\ELfigure{m01-position-vectors}{}

\ELdisp{\vect{A}\ELbr =A_x\uvec{x}+A_y\uvec{y}+A_z\uvec{z},\qquad\ELbr  \vect{B}\ELbr =B_x\uvec{x}+B_y\uvec{y}+B_z\uvec{z}}

\end{ELunit}

\begin{ELjoin}

\begin{itemize}
\tightlist
\item
  Dot product of the two vectors \ELim{\vect{A}} and \ELim{\vect{B}}
\end{itemize}

\begin{ELimportant}{}

\ELdisp{\vect{A}\cdot\vect{B}\ELbr =A_xB_x+A_yB_y+A_zB_z}

\end{ELimportant}

\end{ELjoin}

\begin{ELjoin}

\begin{itemize}
\tightlist
\item
  Cross product of the two vectors \ELim{\vect{A}} and \ELim{\vect{B}}
\end{itemize}

\ELdisp{\begin{aligned}\vect{A}\times\vect{B}={}&A_xB_x\uvec{x}\times\uvec{x}+A_xB_y\uvec{x}\times\uvec{y}+A_xB_z\uvec{x}\times\uvec{z}\\&+A_yB_x\uvec{y}\times\uvec{x}+A_yB_y\uvec{y}\times\uvec{y}+A_yB_z\uvec{y}\times\uvec{z}\\&+A_zB_x\uvec{z}\times\uvec{x}+A_zB_y\uvec{z}\times\uvec{y}+A_zB_z\uvec{z}\times\uvec{z}\\={}&(A_yB_z-A_zB_y)\uvec{x}+(A_zB_x-A_xB_z)\uvec{y}+(A_xB_y-A_yB_x)\uvec{z}\end{aligned}}

\begin{ELimportant}{}

\ELdisp{\vect{A}\times\vect{B}\ELbr =\begin{vmatrix}\uvec{x}&\uvec{y}&\uvec{z}\\A_x&A_y&A_z\\B_x&B_y&B_z\end{vmatrix}}

\end{ELimportant}

\end{ELjoin}

\begin{ELunit}

\begin{itemize}
\tightlist
\item
  Differential length \ELim{\ELto} \textbf{vector}
\end{itemize}

\ELdisp{\dif\vect{\ell}\ELbr =\uvec{x}\,dx+\uvec{y}\,dy+\uvec{z}\,dz}

\begin{itemize}
\tightlist
\item
  Differential surface \ELim{\ELto} \textbf{vector}
\end{itemize}

\ELdisp{d\vect{S}\ELbr =dS\,\uvec{n}}
\ELdisp{d\vect{S}\ELbr =dy\,dz\,\uvec{x},\qquad\ELbr  d\vect{S}\ELbr =dx\,dz\,\uvec{y},\qquad\ELbr  d\vect{S}\ELbr =dx\,dy\,\uvec{z}}

\begin{itemize}
\tightlist
\item
  Differential volume \ELim{\ELto} \textbf{scalar}
\end{itemize}

\ELdisp{dv\ELbr =dx\,dy\,dz}

\ELfigure{m01-cartesian-differentials}{}

\end{ELunit}

\ELnextnum{1-6}

\begin{ELexample}{}

\begin{enumerate}
\def\labelenumi{(\alph{enumi})}
\item
  Find the vector whose initial point is \ELim{P_1(1,3,2)} and whose terminal point is \ELim{P_2(3,-2,4)}.
\item
  What is the length of this vector?
\end{enumerate}

\begin{ELsolution}

\ELfigure{m01-ex6}{}

\begin{enumerate}
\def\labelenumi{(\alph{enumi})}
\tightlist
\item
  \ELdisp{\begin{aligned}\overrightarrow{P_1P_2}&=\overrightarrow{OP_2}-\overrightarrow{OP_1}\\&=(\uvec{x}3-\uvec{y}2+\uvec{z}4)-(\uvec{x}+\uvec{y}3+\uvec{z}2)\\&=\uvec{x}2-\uvec{y}5+\uvec{z}2\end{aligned}}
\item
  \ELdisp{P_1P_2\ELbr =\abs{\overrightarrow{P_1P_2}}\ELbr =\sqrt{2^2+(-5)^2+2^2}\ELbr =\sqrt{33}}
\end{enumerate}

\end{ELsolution}

\end{ELexample}

\ELnextnum{1-7}

\begin{ELexample}{}

Given \ELim{\vect{A}=\uvec{x}5-\uvec{y}2+\uvec{z}}, find the unit vector \ELim{\vect{B}} such that

\begin{itemize}
\item
  \begin{enumerate}
  \def\labelenumi{(\alph{enumi})}
  \tightlist
  \item
    \ELim{\vect{B}\parallel\vect{A}}
  \end{enumerate}
\item
  \begin{enumerate}
  \def\labelenumi{(\alph{enumi})}
  \setcounter{enumi}{1}
  \tightlist
  \item
    \ELim{\vect{B}} lies in the \ELim{xy}-plane and is perpendicular to \ELim{\vect{A}}
  \end{enumerate}
\end{itemize}

\begin{ELsolution}

\begin{enumerate}
\def\labelenumi{(\alph{enumi})}
\tightlist
\item
  \ELdisp{\vect{B}\ELbr =\frac{\vect{A}}{A}\ELbr =\frac{\uvec{x}5-\uvec{y}2+\uvec{z}}{\sqrt{5^2+2^2+1^2}}\ELbr =\frac{\uvec{x}5-\uvec{y}2+\uvec{z}}{\sqrt{30}}}
\item
  \ELdisp{\begin{cases}B=\sqrt{B_x^2+B_y^2+B_z^2}=1\\B_z=0\\\vect{B}\cdot\vect{A}=5B_x-2B_y=0\end{cases}\quad\ELbr \ELbr \Rightarrow\quad\ELbr \begin{cases}B_x=\dfrac{2}{\sqrt{29}}\\\noalign{\vskip 2mm}B_y=\dfrac{5}{\sqrt{29}}\end{cases}}
\end{enumerate}

\end{ELsolution}

\end{ELexample}

\ELhold

\ELsection{Cylindrical Coordinate System}{Cylindrical Coordinate System}

\begin{ELunit}

\ELfigure{m01-cylindrical-surfaces}{}

\begin{itemize}
\tightlist
\item
  Every point \ELim{P(r_1,\phi_1,z_1)} is the intersection of the cylinder \ELim{r=r_1}, the half-plane \ELim{\phi=\phi_1} and the plane \ELim{z=z_1}.
\item
  Unit vector \ELim{\uvec{r}} \ELim{\ELto} the unit vector normal to the \ELim{r}-constant cylinder, in the positive direction
\item
  Unit vector \ELim{\uvec{\phi}} \ELim{\ELto} the unit vector normal to the \ELim{\phi}-constant half-plane, in the positive direction
\item
  Unit vector \ELim{\uvec{z}} \ELim{\ELto} the unit vector normal to the \ELim{z}-constant plane, in the positive direction
\end{itemize}

\end{ELunit}

\begin{ELremark}{}

The directions of the vectors \ELim{\uvec{r}} and \ELim{\uvec{\phi}} change from point to point in space. This must be taken into account in integration.

\end{ELremark}

\begin{ELunit}

\begin{itemize}
\tightlist
\item
  Cross and dot products of the unit vectors
\end{itemize}

\ELdisp{\uvec{r}\times\uvec{\phi}\ELbr =\uvec{z},\qquad\ELbr  \uvec{\phi}\times\uvec{z}\ELbr =\uvec{r},\qquad\ELbr  \uvec{z}\times\uvec{r}\ELbr =\uvec{\phi}}

\ELfigure{m01-cylindrical-cycle}{}

\ELdisp{\uvec{r}\cdot\uvec{r}\ELbr =\uvec{\phi}\cdot\uvec{\phi}\ELbr =\uvec{z}\cdot\uvec{z}\ELbr =1}
\ELdisp{\uvec{r}\cdot\uvec{\phi}\ELbr =\uvec{\phi}\cdot\uvec{z}\ELbr =\uvec{z}\cdot\uvec{r}\ELbr =0}

\end{ELunit}

\begin{ELjoin}

\begin{itemize}
\tightlist
\item
  Representation of an arbitrary vector \ELim{\vect{A}} in cylindrical coordinates
\end{itemize}

\begin{ELimportant}{}

\ELdisp{\vect{A}\ELbr =\uvec{r}A_r+\uvec{\phi}A_\phi+\uvec{z}A_z}
\ELdisp{A\ELbr =\sqrt{A_r^2+A_\phi^2+A_z^2}}

\end{ELimportant}

\end{ELjoin}

\begin{ELunit}

\begin{itemize}
\tightlist
\item
  Position vector
\end{itemize}

\ELfigure{m01-cylindrical-position}{}

\ELdisp{\vect{R}\ELbr =r\uvec{r}+z\uvec{z}}

\end{ELunit}

\ELnextnum{1-8}

\begin{ELexample}{}

The cylindrical coordinates of an arbitrary point \ELim{P} are \ELim{(r,\phi,0)}. Find the unit vector from the point \ELim{z=h} on the \ELim{z}-axis toward the point \ELim{P}.

\begin{ELsolution}

\ELfigure{m01-ex8}{}

\ELdisp{\overrightarrow{QP}\ELbr =\overrightarrow{OP}-\overrightarrow{OQ}\ELbr =(\uvec{r}r)-(\uvec{z}h)}
\ELdisp{\uvec{QP}\ELbr =\frac{\overrightarrow{QP}}{\abs{\overrightarrow{QP}}}\ELbr =\frac{1}{\sqrt{r^2+h^2}}(\uvec{r}r-\uvec{z}h)}

\end{ELsolution}

\end{ELexample}

\begin{ELunit}

\begin{itemize}
\tightlist
\item
  Differential length \ELim{\ELto} \textbf{vector}
\end{itemize}

\ELdisp{\dif\vect{\ell}\ELbr =\uvec{r}\,dr+\uvec{\phi}\,r\,d\phi+\uvec{z}\,dz}

\begin{itemize}
\tightlist
\item
  Differential surface \ELim{\ELto} \textbf{vector}
\end{itemize}

\ELdisp{d\vect{S}\ELbr =r\,d\phi\,dz\,\uvec{r},\qquad\ELbr  d\vect{S}\ELbr =dr\,dz\,\uvec{\phi},\qquad\ELbr  d\vect{S}\ELbr =r\,dr\,d\phi\,\uvec{z}}

\begin{itemize}
\tightlist
\item
  Differential volume \ELim{\ELto} \textbf{scalar}
\end{itemize}

\ELdisp{dv\ELbr =r\,dr\,d\phi\,dz}

\ELfigure{m01-cylindrical-differentials}{}

\begin{itemize}
\tightlist
\item
  Transformation of the representation of a vector from cylindrical to Cartesian coordinates
\end{itemize}

\ELdisp{\vect{A}\ELbr =\uvec{r}A_r+\uvec{\phi}A_\phi+\uvec{z}A_z\quad\ELbr \ELbr \Longrightarrow\quad\ELbr \vect{A}\ELbr =\uvec{x}A_x+\uvec{y}A_y+\uvec{z}A_z}

\ELfigure{m01-cyl-cart-units}{}

\ELdisp{\uvec{r}\ELbr =\cos\phi\,\uvec{x}+\sin\phi\,\uvec{y},\qquad\ELbr  \uvec{\phi}\ELbr =-\sin\phi\,\uvec{x}+\cos\phi\,\uvec{y},\qquad\ELbr  \uvec{z}\ELbr =\uvec{z}}

\ELdisp{\begin{aligned}\vect{A}&=(\cos\phi\,\uvec{x}+\sin\phi\,\uvec{y})A_r\\&\quad+(-\sin\phi\,\uvec{x}+\cos\phi\,\uvec{y})A_\phi\\&\quad+\uvec{z}A_z\end{aligned}\quad\ELbr \ELbr \Longrightarrow\quad\ELbr \begin{aligned}\vect{A}&=\uvec{x}(\cos\phi\,A_r-\sin\phi\,A_\phi)\\&\quad+\uvec{y}(\sin\phi\,A_r+\cos\phi\,A_\phi)\\&\quad+\uvec{z}A_z\end{aligned}}

\begin{itemize}
\tightlist
\item
  Transformation of the representation of a vector from Cartesian to cylindrical coordinates
\end{itemize}

\ELdisp{\vect{A}\ELbr =\uvec{x}A_x+\uvec{y}A_y+\uvec{z}A_z\quad\ELbr \ELbr \Longrightarrow\quad\ELbr \vect{A}\ELbr =\uvec{r}A_r+\uvec{\phi}A_\phi+\uvec{z}A_z}

\ELdisp{\uvec{x}\ELbr =\cos\phi\,\uvec{r}-\sin\phi\,\uvec{\phi},\qquad\ELbr  \uvec{y}\ELbr =\sin\phi\,\uvec{r}+\cos\phi\,\uvec{\phi},\qquad\ELbr  \uvec{z}\ELbr =\uvec{z}}

\ELdisp{\begin{aligned}\vect{A}&=(\cos\phi\,\uvec{r}-\sin\phi\,\uvec{\phi})A_x\\&\quad+(\sin\phi\,\uvec{r}+\cos\phi\,\uvec{\phi})A_y\\&\quad+\uvec{z}A_z\end{aligned}\quad\ELbr \ELbr \Longrightarrow\quad\ELbr \begin{aligned}\vect{A}&=\uvec{r}(\cos\phi\,A_x+\sin\phi\,A_y)\\&\quad+\uvec{\phi}(-\sin\phi\,A_x+\cos\phi\,A_y)\\&\quad+\uvec{z}A_z\end{aligned}}

\end{ELunit}

\begin{ELjoin}

\begin{itemize}
\tightlist
\item
  Transformation of the representation of a vector from cylindrical to Cartesian coordinates
\end{itemize}

\begin{ELimportant}{}

\ELdisp{\begin{bmatrix}A_x\\A_y\\A_z\end{bmatrix}\ELbr =\begin{bmatrix}\cos\phi&-\sin\phi&0\\\sin\phi&\cos\phi&0\\0&0&1\end{bmatrix}\begin{bmatrix}A_r\\A_\phi\\A_z\end{bmatrix}}

\end{ELimportant}

\end{ELjoin}

\begin{ELjoin}

\begin{itemize}
\tightlist
\item
  Transformation of the representation of a vector from Cartesian to cylindrical coordinates
\end{itemize}

\begin{ELimportant}{}

\ELdisp{\begin{bmatrix}A_r\\A_\phi\\A_z\end{bmatrix}\ELbr =\begin{bmatrix}\cos\phi&\sin\phi&0\\-\sin\phi&\cos\phi&0\\0&0&1\end{bmatrix}\begin{bmatrix}A_x\\A_y\\A_z\end{bmatrix}}

\end{ELimportant}

\end{ELjoin}

\begin{ELunit}

\begin{itemize}
\tightlist
\item
  Transformation from cylindrical to Cartesian coordinates
\end{itemize}

\ELdisp{x\ELbr =r\cos\phi,\qquad\ELbr  y\ELbr =r\sin\phi,\qquad\ELbr  z\ELbr =z}

\begin{itemize}
\tightlist
\item
  Transformation from Cartesian to cylindrical coordinates
\end{itemize}

\ELdisp{r\ELbr =\sqrt{x^2+y^2},\qquad\ELbr  \phi\ELbr =\tan^{-1}\frac{y}{x},\qquad\ELbr  z\ELbr =z}

\ELfigure{m01-cyl-cart-coords}{}

\end{ELunit}

\ELhold

\ELsection{Spherical Coordinate System}{Spherical Coordinate System}

\begin{ELunit}

\ELfigure{m01-spherical-surfaces}{}

\begin{itemize}
\tightlist
\item
  Every point \ELim{P(R_1,\theta_1,\phi_1)} is the intersection of the sphere \ELim{R=R_1}, the cone \ELim{\theta=\theta_1} and the half-plane \ELim{\phi=\phi_1}
\item
  Unit vector \ELim{\uvec{R}} \ELim{\ELto} the unit vector normal to the \ELim{R}-constant sphere, in the positive direction
\item
  Unit vector \ELim{\uvec{\theta}} \ELim{\ELto} the unit vector normal to the \ELim{\theta}-constant cone, in the positive direction
\item
  Unit vector \ELim{\uvec{\phi}} \ELim{\ELto} the unit vector normal to the \ELim{\phi}-constant half-plane, in the positive direction
\end{itemize}

\end{ELunit}

\begin{ELremark}{}

The directions of the vectors \ELim{\uvec{R}}, \ELim{\uvec{\theta}} and \ELim{\uvec{\phi}} change from point to point in space. This must be taken into account in integration.

\end{ELremark}

\begin{ELunit}

\begin{itemize}
\tightlist
\item
  Cross and dot products of the unit vectors
\end{itemize}

\ELdisp{\uvec{R}\times\uvec{\theta}\ELbr =\uvec{\phi},\qquad\ELbr  \uvec{\theta}\times\uvec{\phi}\ELbr =\uvec{R},\qquad\ELbr  \uvec{\phi}\times\uvec{R}\ELbr =\uvec{\theta}}

\ELfigure{m01-spherical-cycle}{}

\ELdisp{\uvec{R}\cdot\uvec{R}\ELbr =\uvec{\theta}\cdot\uvec{\theta}\ELbr =\uvec{\phi}\cdot\uvec{\phi}\ELbr =1}
\ELdisp{\uvec{R}\cdot\uvec{\theta}\ELbr =\uvec{\theta}\cdot\uvec{\phi}\ELbr =\uvec{\phi}\cdot\uvec{R}\ELbr =0}

\end{ELunit}

\begin{ELjoin}

\begin{itemize}
\tightlist
\item
  Representation of an arbitrary vector \ELim{\vect{A}} in spherical coordinates
\end{itemize}

\begin{ELimportant}{}

\ELdisp{\vect{A}\ELbr =\uvec{R}A_R+\uvec{\theta}A_\theta+\uvec{\phi}A_\phi}
\ELdisp{A\ELbr =\sqrt{A_R^2+A_\theta^2+A_\phi^2}}

\end{ELimportant}

\end{ELjoin}

\begin{ELunit}

\begin{itemize}
\tightlist
\item
  Position vector
\end{itemize}

\ELfigure{m01-spherical-position}{}

\ELdisp{\vect{R}\ELbr =R_1\uvec{R}}

\begin{itemize}
\tightlist
\item
  Differential length \ELim{\ELto} \textbf{vector}
\end{itemize}

\ELdisp{\dif\vect{\ell}\ELbr =\uvec{R}\,dR+\uvec{\theta}\,R\,d\theta+\uvec{\phi}\,R\sin\theta\,d\phi}

\begin{itemize}
\tightlist
\item
  Differential surface \ELim{\ELto} \textbf{vector}
\end{itemize}

\ELdisp{d\vect{S}\ELbr =R^2\sin\theta\,d\theta\,d\phi\,\uvec{R},\qquad\ELbr  d\vect{S}\ELbr =R\sin\theta\,dR\,d\phi\,\uvec{\theta},\qquad\ELbr  d\vect{S}\ELbr =R\,dR\,d\theta\,\uvec{\phi}}

\begin{itemize}
\tightlist
\item
  Differential volume \ELim{\ELto} \textbf{scalar}
\end{itemize}

\ELdisp{dv\ELbr =R^2\sin\theta\,dR\,d\theta\,d\phi}

\ELfigure{m01-spherical-differentials}{}

\begin{itemize}
\tightlist
\item
  Transformation of the representation of a vector from spherical to cylindrical coordinates
\end{itemize}

\ELdisp{\vect{A}\ELbr =\uvec{R}A_R+\uvec{\theta}A_\theta+\uvec{\phi}A_\phi\quad\ELbr \ELbr \Longrightarrow\quad\ELbr \vect{A}\ELbr =\uvec{r}A_r+\uvec{\phi}A_\phi+\uvec{z}A_z}

\ELfigure{m01-sph-cyl-units}{}

\ELdisp{\uvec{R}\ELbr =\sin\theta\,\uvec{r}+\cos\theta\,\uvec{z},\qquad\ELbr  \uvec{\theta}\ELbr =\cos\theta\,\uvec{r}-\sin\theta\,\uvec{z},\qquad\ELbr  \uvec{\phi}\ELbr =\uvec{\phi}}

\ELdisp{\begin{aligned}\vect{A}&=(\sin\theta\,\uvec{r}+\cos\theta\,\uvec{z})A_R\\&\quad+(\cos\theta\,\uvec{r}-\sin\theta\,\uvec{z})A_\theta\\&\quad+\uvec{\phi}A_\phi\end{aligned}\quad\ELbr \ELbr \Longrightarrow\quad\ELbr \begin{aligned}\vect{A}&=\uvec{r}(\sin\theta\,A_R+\cos\theta\,A_\theta)\\&\quad+\uvec{\phi}A_\phi\\&\quad+\uvec{z}(\cos\theta\,A_R-\sin\theta\,A_\theta)\end{aligned}}

\begin{itemize}
\tightlist
\item
  Transformation of the representation of a vector from cylindrical to spherical coordinates
\end{itemize}

\ELdisp{\vect{A}\ELbr =\uvec{r}A_r+\uvec{\phi}A_\phi+\uvec{z}A_z\quad\ELbr \ELbr \Longrightarrow\quad\ELbr \vect{A}\ELbr =\uvec{R}A_R+\uvec{\theta}A_\theta+\uvec{\phi}A_\phi}

\ELdisp{\uvec{r}\ELbr =\sin\theta\,\uvec{R}+\cos\theta\,\uvec{\theta},\qquad\ELbr  \uvec{\phi}\ELbr =\uvec{\phi},\qquad\ELbr  \uvec{z}\ELbr =\cos\theta\,\uvec{R}-\sin\theta\,\uvec{\theta}}

\ELdisp{\begin{aligned}\vect{A}&=(\sin\theta\,\uvec{R}+\cos\theta\,\uvec{\theta})A_r\\&\quad+\uvec{\phi}A_\phi\\&\quad+(\cos\theta\,\uvec{R}-\sin\theta\,\uvec{\theta})A_z\end{aligned}\quad\ELbr \ELbr \Longrightarrow\quad\ELbr \begin{aligned}\vect{A}&=\uvec{R}(\sin\theta\,A_r+\cos\theta\,A_z)\\&\quad+\uvec{\theta}(\cos\theta\,A_r-\sin\theta\,A_z)\\&\quad+\uvec{\phi}A_\phi\end{aligned}}

\end{ELunit}

\begin{ELjoin}

\begin{itemize}
\tightlist
\item
  Transformation of the representation of a vector from spherical to cylindrical coordinates
\end{itemize}

\begin{ELimportant}{}

\ELdisp{\begin{bmatrix}A_r\\A_\phi\\A_z\end{bmatrix}\ELbr =\begin{bmatrix}\sin\theta&\cos\theta&0\\0&0&1\\\cos\theta&-\sin\theta&0\end{bmatrix}\begin{bmatrix}A_R\\A_\theta\\A_\phi\end{bmatrix}}

\end{ELimportant}

\end{ELjoin}

\begin{ELjoin}

\begin{itemize}
\tightlist
\item
  Transformation of the representation of a vector from cylindrical to spherical coordinates
\end{itemize}

\begin{ELimportant}{}

\ELdisp{\begin{bmatrix}A_R\\A_\theta\\A_\phi\end{bmatrix}\ELbr =\begin{bmatrix}\sin\theta&0&\cos\theta\\\cos\theta&0&-\sin\theta\\0&1&0\end{bmatrix}\begin{bmatrix}A_r\\A_\phi\\A_z\end{bmatrix}}

\end{ELimportant}

\end{ELjoin}

\begin{ELunit}

\begin{itemize}
\tightlist
\item
  Transformation of the representation of a vector from Cartesian to spherical coordinates
\end{itemize}

\ELdisp{\vect{A}\ELbr =\uvec{x}A_x+\uvec{y}A_y+\uvec{z}A_z\quad\ELbr \ELbr \Longrightarrow\quad\ELbr \vect{A}\ELbr =\uvec{R}A_R+\uvec{\theta}A_\theta+\uvec{\phi}A_\phi}

\ELdisp{\left\{\begin{aligned}\uvec{x}&=\cos\phi\,\uvec{r}-\sin\phi\,\uvec{\phi}\\\uvec{y}&=\sin\phi\,\uvec{r}+\cos\phi\,\uvec{\phi}\\\uvec{z}&=\uvec{z}\end{aligned}\right.\qquad\ELbr \left\{\begin{aligned}\uvec{r}&=\sin\theta\,\uvec{R}+\cos\theta\,\uvec{\theta}\\\uvec{\phi}&=\uvec{\phi}\\\uvec{z}&=\cos\theta\,\uvec{R}-\sin\theta\,\uvec{\theta}\end{aligned}\right.}

\ELdisp{\Longrightarrow\quad\ELbr \left\{\begin{aligned}\uvec{x}&=\sin\theta\cos\phi\,\uvec{R}+\cos\theta\cos\phi\,\uvec{\theta}-\sin\phi\,\uvec{\phi}\\\uvec{y}&=\sin\theta\sin\phi\,\uvec{R}+\cos\theta\sin\phi\,\uvec{\theta}+\cos\phi\,\uvec{\phi}\\\uvec{z}&=\cos\theta\,\uvec{R}-\sin\theta\,\uvec{\theta}\end{aligned}\right.}

\end{ELunit}

\begin{ELimportant}{}

\ELdisp{\begin{bmatrix}A_R\\A_\theta\\A_\phi\end{bmatrix}\ELbr =\begin{bmatrix}\sin\theta\cos\phi&\sin\theta\sin\phi&\cos\theta\\\cos\theta\cos\phi&\cos\theta\sin\phi&-\sin\theta\\-\sin\phi&\cos\phi&0\end{bmatrix}\begin{bmatrix}A_x\\A_y\\A_z\end{bmatrix}}

\end{ELimportant}

\begin{ELunit}

\begin{itemize}
\tightlist
\item
  Transformation of the representation of a vector from spherical to Cartesian coordinates
\end{itemize}

\ELdisp{\vect{A}\ELbr =\uvec{R}A_R+\uvec{\theta}A_\theta+\uvec{\phi}A_\phi\quad\ELbr \ELbr \Longrightarrow\quad\ELbr \vect{A}\ELbr =\uvec{x}A_x+\uvec{y}A_y+\uvec{z}A_z}

\ELdisp{\left\{\begin{aligned}\uvec{R}&=\sin\theta\,\uvec{r}+\cos\theta\,\uvec{z}\\\uvec{\theta}&=\cos\theta\,\uvec{r}-\sin\theta\,\uvec{z}\\\uvec{\phi}&=\uvec{\phi}\end{aligned}\right.\qquad\ELbr \left\{\begin{aligned}\uvec{r}&=\cos\phi\,\uvec{x}+\sin\phi\,\uvec{y}\\\uvec{\phi}&=-\sin\phi\,\uvec{x}+\cos\phi\,\uvec{y}\\\uvec{z}&=\uvec{z}\end{aligned}\right.}

\ELdisp{\Longrightarrow\quad\ELbr \left\{\begin{aligned}\uvec{R}&=\sin\theta\cos\phi\,\uvec{x}+\sin\theta\sin\phi\,\uvec{y}+\cos\theta\,\uvec{z}\\\uvec{\theta}&=\cos\theta\cos\phi\,\uvec{x}+\cos\theta\sin\phi\,\uvec{y}-\sin\theta\,\uvec{z}\\\uvec{\phi}&=-\sin\phi\,\uvec{x}+\cos\phi\,\uvec{y}\end{aligned}\right.}

\end{ELunit}

\begin{ELimportant}{}

\ELdisp{\begin{bmatrix}A_x\\A_y\\A_z\end{bmatrix}\ELbr =\begin{bmatrix}\sin\theta\cos\phi&\cos\theta\cos\phi&-\sin\phi\\\sin\theta\sin\phi&\cos\theta\sin\phi&\cos\phi\\\cos\theta&-\sin\theta&0\end{bmatrix}\begin{bmatrix}A_R\\A_\theta\\A_\phi\end{bmatrix}}

\end{ELimportant}

\begin{ELunit}

\begin{itemize}
\tightlist
\item
  Transformation from cylindrical to spherical coordinates
\end{itemize}

\ELdisp{R\ELbr =\sqrt{r^2+z^2},\qquad\ELbr  \theta\ELbr =\tan^{-1}\frac{r}{z},\qquad\ELbr  \phi\ELbr =\phi}

\begin{itemize}
\tightlist
\item
  Transformation from spherical to cylindrical coordinates
\end{itemize}

\ELdisp{r\ELbr =R\sin\theta,\qquad\ELbr  \phi\ELbr =\phi,\qquad\ELbr  z\ELbr =R\cos\theta}

\begin{itemize}
\tightlist
\item
  Transformation from spherical to Cartesian coordinates
\end{itemize}

\ELdisp{x\ELbr =R\sin\theta\cos\phi,\qquad\ELbr  y\ELbr =R\sin\theta\sin\phi,\qquad\ELbr  z\ELbr =R\cos\theta}

\begin{itemize}
\tightlist
\item
  Transformation from Cartesian to spherical coordinates
\end{itemize}

\ELdisp{R\ELbr =\sqrt{x^2+y^2+z^2},\qquad\ELbr  \theta\ELbr =\tan^{-1}\frac{\sqrt{x^2+y^2}}{z},\qquad\ELbr  \phi\ELbr =\tan^{-1}\frac{y}{x}}

\ELfigure{m01-sph-coords}{}

\end{ELunit}

\ELhold

\ELsection{Integrals Containing Vector Functions}{Integrals Containing Vector Functions}

\begin{ELunit}

\ELfigure{m01-vector-analysis-map-calc}{}

\begin{itemize}
\tightlist
\item
  The kinds of integrals containing vector functions that we deal with
\end{itemize}

\ELdisp{\int_C\vect{F}\,d\ell\qquad\ELbr  \iint_S\vect{F}\,ds\qquad\ELbr  \iiint_V\vect{F}\,dv}
\ELdisp{\int_C V\,\dif\vect{\ell}}
\ELdisp{\int_C\vect{F}\cdot\dif\vect{\ell}\qquad\ELbr  \iint_S\vect{A}\cdot d\vect{s}}
\ELdisp{\int_C\vect{F}\times\dif\vect{\ell}}

\end{ELunit}

\begin{ELremark}{}

Important points about integration involving vector functions

\begin{itemize}
\tightlist
\item
  Using an appropriate coordinate system
\item
  Correct determination of the differential length, surface and volume
\item
  Attention to the variation of the unit vectors \ELim{\uvec{r}}, \ELim{\uvec{R}}, \ELim{\uvec{\theta}} and \ELim{\uvec{\phi}} over the range of integration
\end{itemize}

\end{ELremark}

\begin{ELunit}

\begin{itemize}
\tightlist
\item
  The integral of a vector function over a specified curve, surface and volume
\end{itemize}

\ELdisp{\int_C\vect{F}\,d\ell\qquad\ELbr  \iint_S\vect{F}\,ds\qquad\ELbr  \iiint_V\vect{F}\,dv}

\begin{itemize}
\tightlist
\item
  The result is a vector.
\item
  Important applications in electromagnetics: computing the electric field intensity due to a line, surface and volume charge distribution
\end{itemize}

\ELdisp{\left\{\begin{aligned}\vect{E}&=\frac{1}{4\pi\varepsilon_0}\int_C\frac{\rho_\ell\,d\ell'(\vect{R}-\vect{R}')}{\abs{\vect{R}-\vect{R}'}^3}\\\vect{E}&=\frac{1}{4\pi\varepsilon_0}\iint_S\frac{\rho_s\,ds'(\vect{R}-\vect{R}')}{\abs{\vect{R}-\vect{R}'}^3}\\\vect{E}&=\frac{1}{4\pi\varepsilon_0}\iiint_V\frac{\rho_v\,dv'(\vect{R}-\vect{R}')}{\abs{\vect{R}-\vect{R}'}^3}\end{aligned}\right.}

\end{ELunit}

\ELnextnum{1-9}

\begin{ELexample}{}

Evaluate the following integral over the curve \ELim{C}.
\ELdisp{\int_C\vect{F}\,d\ell,\qquad\ELbr  \vect{F}\ELbr =3\uvec{r}}

\ELfigure{m01-ex9}{}

\begin{ELsolution}

\ELfigure{m01-ex9-solution}{}

\ELdisp{d\ell\ELbr =r\,d\phi\ELbr =a\,d\phi}
\ELdisp{\int_C\vect{F}\,d\ell\ELbr =\int_0^\pi 3\uvec{r}\,a\,d\phi\ELbr =3a\int_0^\pi(\cos\phi\,\uvec{x}+\sin\phi\,\uvec{y})\,d\phi\ELbr =6a\uvec{y}}

\end{ELsolution}

\end{ELexample}

\begin{ELunit}

\begin{itemize}
\tightlist
\item
  The integral of a scalar function along a specified path
\end{itemize}

\ELdisp{\int_C V\,\dif\vect{\ell}}

\begin{itemize}
\tightlist
\item
  The result is a vector.
\item
  Important applications in electromagnetics: computing the magnetic potential due to a line current
\end{itemize}

\ELdisp{\vect{A}\ELbr =\frac{\mu_0I}{4\pi}\oint_C\frac{\dif\vect{\ell}'}{\abs{\vect{R}-\vect{R}'}}}

\end{ELunit}

\ELnextnum{1-10}

\begin{ELexample}{}

Evaluate the following integral along the following paths

\ELdisp{\int_O^P r^2\,\dif\vect{\ell},\qquad\ELbr  r^2\ELbr =x^2+y^2}

\begin{itemize}
\tightlist
\item
  Path \ELim{OP}
\item
  Path \ELim{OP_1P}
\item
  Path \ELim{OP_2P}
\end{itemize}

\ELfigure{m01-ex10}{}

\begin{ELsolution}

\begin{itemize}
\tightlist
\item
  Along the path \ELim{OP}
\end{itemize}

\ELdisp{\begin{aligned}\int_O^P r^2\,\dif\vect{\ell}&=\uvec{r}\int_0^{\sqrt2}r^2\,dr=\uvec{r}\,\frac{2\sqrt2}{3}\\&=\frac{2\sqrt2}{3}(\uvec{x}\cos45^\circ+\uvec{y}\sin45^\circ)\\&=\uvec{x}\frac23+\uvec{y}\frac23\end{aligned}}

\begin{itemize}
\tightlist
\item
  Along the path \ELim{OP_1P}
\end{itemize}

\ELdisp{\begin{aligned}\int_O^P(x^2+y^2)\,\dif\vect{\ell}&=\uvec{y}\int_O^{P_1}y^2\,dy+\uvec{x}\int_{P_1}^P(x^2+1)\,dx\\&=\uvec{y}\tfrac13y^3\Big|_0^1+\uvec{x}\left(\tfrac13x^3+x\right)\Big|_0^1\\&=\uvec{x}\frac43+\uvec{y}\frac13\end{aligned}}

\begin{itemize}
\tightlist
\item
  Along the path \ELim{OP_2P}
\end{itemize}

\ELdisp{\begin{aligned}\int_O^P(x^2+y^2)\,\dif\vect{\ell}&=\uvec{x}\int_O^{P_2}x^2\,dx+\uvec{y}\int_{P_2}^P(1+y^2)\,dy\\&=\uvec{x}\tfrac13x^3\Big|_0^1+\uvec{y}\left(y+\tfrac13y^3\right)\Big|_0^1\\&=\uvec{x}\frac13+\uvec{y}\frac43\end{aligned}}

\end{ELsolution}

\end{ELexample}

\begin{ELunit}

\begin{itemize}
\tightlist
\item
  Scalar line integral
\end{itemize}

\ELdisp{\int_C\vect{F}\cdot\dif\vect{\ell}}

\begin{itemize}
\tightlist
\item
  The result is a scalar.
\item
  If \ELim{\vect{F}} is a force

  \begin{itemize}
  \tightlist
  \item
    the work done by the force in moving a body from point \ELim{P_1} to point \ELim{P_2} along the specified path \ELim{C}
  \end{itemize}
\item
  Important applications in electromagnetics

  \begin{itemize}
  \tightlist
  \item
    Computing the electric potential difference between the points \ELim{P_1} and \ELim{P_2}:
    \ELdisp{V\ELbr =-\int_{P_1}^{P_2}\vect{E}\cdot\dif\vect{\ell}}
  \item
    Ampère's circuital law:
    \ELdisp{\mu_0I\ELbr =-\oint_C\vect{B}\cdot\dif\vect{\ell}}
  \end{itemize}
\end{itemize}

\end{ELunit}

\ELnextnum{1-11}

\begin{ELexample}{}

Evaluate the following integral along the quarter-circle path of the figure.

\ELdisp{\int_A^B\vect{F}\cdot\dif\vect{\ell},\qquad\ELbr  \vect{F}\ELbr =\uvec{x}xy-\uvec{y}2x}

\ELfigure{m01-ex11}{}

\begin{ELsolution}

\ELdisp{\begin{bmatrix}F_r\\F_\phi\\F_z\end{bmatrix}\ELbr =\begin{bmatrix}\cos\phi&\sin\phi&0\\-\sin\phi&\cos\phi&0\\0&0&1\end{bmatrix}\begin{bmatrix}xy\\-2x\\0\end{bmatrix}}
\ELdisp{\vect{F}\ELbr =\uvec{r}(xy\cos\phi-2x\sin\phi)-\uvec{\phi}(xy\sin\phi+2x\cos\phi)}
\ELdisp{\dif\vect{\ell}\ELbr =\uvec{\phi}\,3\,d\phi\qquad\ELbr  \vect{F}\cdot\dif\vect{\ell}\ELbr =-3(xy\sin\phi+2x\cos\phi)\,d\phi}
\ELdisp{x\ELbr =3\cos\phi\qquad\ELbr  y\ELbr =3\sin\phi}
\ELdisp{\int_A^B\vect{F}\cdot\dif\vect{\ell}\ELbr =\int_0^{\pi/2}-3(9\sin^2\phi\cos\phi+6\cos^2\phi)\,d\phi\ELbr =-9\left(1+\frac{\pi}{2}\right)}

\end{ELsolution}

\end{ELexample}

\begin{ELunit}

\begin{itemize}
\tightlist
\item
  Scalar surface integral
\end{itemize}

\ELdisp{\iint_S\vect{A}\cdot d\vect{s}}

\begin{itemize}
\tightlist
\item
  The result is a scalar.
\item
  Result of the integral \ELim{\ELto} the flux of the vector field \ELim{\vect{A}} passing through the surface \ELim{S}
\item
  Vector differential surface
\end{itemize}

\ELdisp{d\vect{s}\ELbr =\uvec{n}\,ds}

\ELfigure{m01-surface-normals}{Direction of \ELim{\mathbf{a}_n} for a closed surface: normal to the
surface, pointing out of the volume; for an open surface: depends on the
direction of travel along the rim of the open surface, by the right-hand
rule}

\end{ELunit}

\begin{ELjoin}

\begin{itemize}
\tightlist
\item
  Important applications in electromagnetics

  \begin{itemize}
  \tightlist
  \item
    Gauss's law in free space:
    \ELdisp{\oiint_S\vect{E}\cdot d\vect{s}\ELbr =\frac{Q}{\varepsilon_0}}
  \item
    Computing the electric current:
    \ELdisp{\iint_S\vect{J}\cdot d\vect{s}\ELbr =I}
  \end{itemize}
\end{itemize}

\ELnextnum{1-12}

\begin{ELexample}{}

Evaluate the following integral over the closed surface about the \ELim{z}-axis specified by \ELim{z=\pm3} and \ELim{r=2}

\ELdisp{\oint_S\vect{F}\cdot d\vect{s},\qquad\ELbr  \vect{F}\ELbr =\uvec{r}\frac{k_1}{r}+\uvec{z}k_2z}

\ELfigure{m01-ex12}{}

\begin{ELsolution}

\ELdisp{\oint_S\vect{F}\cdot d\vect{s}\ELbr =\oint_S\vect{F}\cdot\uvec{n}\,ds\ELbr =\int_{\substack{\text{top}\\\text{face}}}\vect{F}\cdot\uvec{n}\,ds+\int_{\substack{\text{bottom}\\\text{face}}}\vect{F}\cdot\uvec{n}\,ds+\int_{\substack{\text{side}\\\text{wall}}}\vect{F}\cdot\uvec{n}\,ds}

\begin{itemize}
\tightlist
\item
  Top face
\end{itemize}

\ELdisp{z\ELbr =3,\quad\ELbr  \uvec{n}\ELbr =\uvec{z},\quad\ELbr  \vect{F}\cdot\uvec{n}\ELbr =k_2z\ELbr =3k_2,\quad\ELbr  ds\ELbr =r\,dr\,d\phi}
\ELdisp{\int_{\substack{\text{top}\\\text{face}}}\vect{F}\cdot\uvec{n}\,ds\ELbr =\int_0^{2\pi}\!\!\int_0^2 3k_2\,r\,dr\,d\phi\ELbr =12\pi k_2}

\begin{itemize}
\tightlist
\item
  Bottom face
\end{itemize}

\ELdisp{z\ELbr =-3,\quad\ELbr  \uvec{n}\ELbr =-\uvec{z},\quad\ELbr  \vect{F}\cdot\uvec{n}\ELbr =-k_2z\ELbr =3k_2,\quad\ELbr  ds\ELbr =r\,dr\,d\phi}
\ELdisp{\int_{\substack{\text{bottom}\\\text{face}}}\vect{F}\cdot\uvec{n}\,ds\ELbr =12\pi k_2}

\begin{itemize}
\tightlist
\item
  Side wall
\end{itemize}

\ELdisp{r\ELbr =2,\quad\ELbr  \uvec{n}\ELbr =\uvec{r},\quad\ELbr  \vect{F}\cdot\uvec{n}\ELbr =\frac{k_1}{r}\ELbr =\frac{k_1}{2},\quad\ELbr  ds\ELbr =r\,d\phi\,dz\ELbr =2\,d\phi\,dz}
\ELdisp{\int_{\substack{\text{side}\\\text{wall}}}\vect{F}\cdot\uvec{n}\,ds\ELbr =\int_{-3}^3\!\int_0^{2\pi}k_1\,d\phi\,dz\ELbr =12\pi k_1}

\ELdisp{\oint_S\vect{F}\cdot d\vect{s}\ELbr =12\pi k_2+12\pi k_2+12\pi k_1\ELbr =12\pi(k_1+2k_2)}

\end{ELsolution}

\end{ELexample}

\end{ELjoin}

\begin{ELjoin}

\begin{itemize}
\tightlist
\item
  Vector line integral
\end{itemize}

\ELdisp{\int_C\vect{F}\times\dif\vect{\ell}}

\begin{itemize}
\tightlist
\item
  The result is a vector.
\item
  Important applications in electromagnetics

  \begin{itemize}
  \tightlist
  \item
    Biot--Savart law:
    \ELdisp{\vect{B}\ELbr =\frac{\mu_0I}{4\pi}\oint_C\frac{\dif\vect{\ell}\times(\vect{R}-\vect{R}')}{\abs{\vect{R}-\vect{R}'}^3}}
  \item
    Computing the magnetic force on a current-carrying wire in a magnetic field:
    \ELdisp{\vect{F}\ELbr =I\oint_C\dif\vect{\ell}\times\vect{B}}
  \end{itemize}
\end{itemize}

\ELnextnum{1-13}

\begin{ELexample}{}

Evaluate the following integral along the path \ELim{C}.

\ELdisp{\int_C\vect{F}\times\dif\vect{\ell},\qquad\ELbr  \vect{F}\ELbr =2\uvec{z}}

\ELfigure{m01-ex13}{}

\begin{ELsolution}

\ELdisp{\dif\vect{\ell}\ELbr =r\,d\phi\,\uvec{\phi}\ELbr =a\,d\phi\,\uvec{\phi}}
\ELdisp{\begin{aligned}\int_C\vect{F}\times\dif\vect{\ell}&=\int_{2\pi}^0 2\uvec{z}\times a\,d\phi\,\uvec{\phi}=-2a\int_{2\pi}^0\uvec{r}\,d\phi\\&=-2a\int_{2\pi}^0(\cos\phi\,\uvec{x}+\sin\phi\,\uvec{y})\,d\phi=0\end{aligned}}

\end{ELsolution}

\end{ELexample}

\end{ELjoin}

\ELhold

\ELsection{Differential Operations - Gradient}{Differential Operations - Gradient}

\begin{ELunit}

\begin{itemize}
\tightlist
\item
  \textbf{Prerequisite concepts for the discussion of the gradient}

  \begin{itemize}
  \tightlist
  \item
    The concept of the derivative of a function of one variable

    \begin{itemize}
    \tightlist
    \item
      The rate of change of the function at the point of interest
    \end{itemize}
  \item
    The concept of the directional derivative of a function of several variables

    \begin{itemize}
    \tightlist
    \item
      The rate of change of the function at the point of interest and in a specified direction
    \end{itemize}
  \end{itemize}
\end{itemize}

\ELfigure{m01-derivative-concepts}{}

\end{ELunit}

\begin{ELdefinition}{Gradient}

The gradient of a \textbf{scalar quantity} is a \textbf{vector quantity} that lies in the direction in which the scalar quantity increases most, and whose magnitude equals the maximum rate of increase of the scalar quantity.

\ELdisp{\operatorname{grad}v\ELbr =\nabla v\ELbr =\uvec{n}\frac{dv}{dn}}

\end{ELdefinition}

\begin{ELunit}

\begin{itemize}
\tightlist
\item
  The directional derivative of \ELim{v} in the direction of \ELim{\dif\vect{\ell}}:
\end{itemize}

\ELdisp{\frac{dv}{d\ell}\ELbr =\nabla v\cdot\uvec{\ell}}

\begin{itemize}
\tightlist
\item
  The gradient of \ELim{v} is normal to the surfaces of constant \ELim{v}.
\end{itemize}

\ELfigure{m01-gradient-map}{}

\end{ELunit}

\begin{ELjoin}

\begin{itemize}
\tightlist
\item
  The gradient in Cartesian coordinates is as follows
\end{itemize}

\begin{ELimportant}{}

\ELdisp{\nabla v\ELbr =\frac{\partial v}{\partial x}\uvec{x}+\frac{\partial v}{\partial y}\uvec{y}+\frac{\partial v}{\partial z}\uvec{z}}

\end{ELimportant}

\end{ELjoin}

\begin{ELjoin}

\begin{itemize}
\tightlist
\item
  The general expression of the gradient
\end{itemize}

\begin{ELimportant}{}

\ELdisp{\nabla V\ELbr =\uvec{u_1}\frac{\partial V}{h_1\,\partial u_1}+\uvec{u_2}\frac{\partial V}{h_2\,\partial u_2}+\uvec{u_3}\frac{\partial V}{h_3\,\partial u_3}}

\end{ELimportant}

\end{ELjoin}

\Needspace{10\baselineskip}
\begin{longtable}[c]{llll}
\toprule\noalign{}
\ELtablehead Metric coefficients & Cartesian & Cylindrical & Spherical \\\noalign{\vskip 4pt}
\midrule\noalign{}
\endhead
\bottomrule\noalign{}
\endlastfoot
\ELim{h_1} & \ELim{1} & \ELim{1} & \ELim{1} \\\noalign{\vskip 4pt}
\ELim{h_2} & \ELim{1} & \ELim{r} & \ELim{R} \\\noalign{\vskip 4pt}
\ELim{h_3} & \ELim{1} & \ELim{1} & \ELim{R\sin\theta} \\\noalign{\vskip 4pt}
\end{longtable}

\begin{ELunit}

\begin{itemize}
\tightlist
\item
  \textbf{The concept of the gradient}

  \begin{itemize}
  \tightlist
  \item
    Representing a hill by contours of height \ELim{\ELto} \ELim{h(x,y)}
  \item
    If a ball is placed at point \ELim{P}, in which direction does it move and what is the magnitude of the force acting on it?

    \begin{itemize}
    \tightlist
    \item
      The answer is proportional to \ELim{-\nabla h}.
    \end{itemize}
  \end{itemize}
\end{itemize}

\ELfigure{m01-gradient-hill}{}

\end{ELunit}

\ELnextnum{1-14}

\begin{ELexample}{}

The electric field intensity \ELim{\vect{E}} can be obtained as the negative of the gradient of the scalar electric potential \ELim{V}. Find \ELim{\vect{E}} at the point \ELim{(1,1,0)} if
\ELdisp{V\ELbr =E_0R\cos\theta}

\begin{ELsolution}

\ELdisp{\vect{E}\ELbr =-\nabla V}
\ELdisp{\begin{aligned}\vect{E}&=-\left[\uvec{R}\frac{\partial}{\partial R}+\uvec{\theta}\frac{\partial}{R\,\partial\theta}+\uvec{\phi}\frac{\partial}{R\sin\theta\,\partial\phi}\right]E_0R\cos\theta\\&=-(\uvec{R}\cos\theta-\uvec{\theta}\sin\theta)E_0\end{aligned}}
\ELdisp{\vect{E}\ELbr =-\uvec{z}E_0}

\end{ELsolution}

\end{ELexample}

\ELhold

\ELsection{Differential Operations - Divergence}{Differential Operations - Divergence}

\begin{ELunit}

\begin{itemize}
\tightlist
\item
  \textbf{Prerequisite concepts for the discussion of the divergence}

  \begin{itemize}
  \tightlist
  \item
    In studying vector fields it is easier to represent the variations of the field graphically by flux lines. These are directed lines or curves that indicate the direction of the vector field at each point.
  \end{itemize}
\end{itemize}

\ELfigure{m01-flux-lines}{}

\begin{itemize}
\tightlist
\item
  The flux of a vector field is analogous to the flow of a fluid such as water.
\item
  If \ELim{\vect{A}} is a flux density vector, the total flux leaving the surface \ELim{S} equals
\end{itemize}

\ELdisp{\oint_S\vect{A}\cdot d\vect{s}}

\begin{itemize}
\tightlist
\item
  In a volume with a closed surface, there will be an excess of inward or outward flux only when the volume contains a sink or a source, respectively.

  \begin{itemize}
  \tightlist
  \item
    The net outward flux from \ELim{v_a} is positive: there is a source in the volume \ELim{v_a}
  \item
    The net outward flux from \ELim{v_b} is negative: there is a sink in the volume \ELim{v_b}
  \item
    The net outward flux from \ELim{v_c} is zero.
  \end{itemize}
\end{itemize}

\ELfigure{m01-source-sink}{}

\end{ELunit}

\begin{ELdefinition}{Divergence}

The divergence of a \textbf{vector field} \ELim{\vect{A}} at a point is the net outward flux of \ELim{\vect{A}} per unit volume as the volume about the point tends to zero

\ELdisp{\operatorname{div}\vect{A}\triangleq\lim_{\Delta v\to0}\frac{\oint_S\vect{A}\cdot d\vect{s}}{\Delta v}}

\end{ELdefinition}

\begin{ELjoin}

\begin{itemize}
\tightlist
\item
  In Cartesian coordinates the divergence is as follows
\end{itemize}

\begin{ELimportant}{}

\ELdisp{\operatorname{div}\vect{A}\ELbr =\frac{\partial A_x}{\partial x}+\frac{\partial A_y}{\partial y}+\frac{\partial A_z}{\partial z}\qquad\ELbr \qquad\ELbr  \nabla\cdot\vect{A}\equiv\operatorname{div}\vect{A}}

\end{ELimportant}

\end{ELjoin}

\begin{ELjoin}

\begin{itemize}
\tightlist
\item
  The general expression of the divergence
\end{itemize}

\begin{ELimportant}{}

\ELdisp{\nabla\cdot\vect{A}\ELbr =\frac{1}{h_1h_2h_3}\left[\frac{\partial}{\partial u_1}(h_2h_3A_1)+\frac{\partial}{\partial u_2}(h_1h_3A_2)+\frac{\partial}{\partial u_3}(h_1h_2A_3)\right]}

\end{ELimportant}

\end{ELjoin}

\ELnextnum{1-15}

\begin{ELexample}{}

The magnetic flux density \ELim{\vect{B}} outside a long current-carrying wire is circular and inversely proportional to the distance from the axis of the wire. Find the divergence of \ELim{\vect{B}}.

\begin{ELsolution}

\begin{itemize}
\tightlist
\item
  Assuming the wire coincides with the \ELim{z}-axis
\end{itemize}

\ELdisp{\vect{B}\ELbr =\uvec{\phi}\frac{k}{r}}
\ELdisp{\nabla\cdot\vect{B}\ELbr =\frac1r\frac{\partial}{\partial r}(rB_r)+\frac1r\frac{\partial B_\phi}{\partial\phi}+\frac{\partial B_z}{\partial z}}
\ELdisp{\nabla\cdot\vect{B}\ELbr =0}

\ELfigure{m01-ex15}{}

\end{ELsolution}

\end{ELexample}

\ELhold

\ELsection{Differential Operations - Curl}{Differential Operations - Curl}

\begin{ELunit}

\begin{itemize}
\tightlist
\item
  \textbf{Prerequisite concepts for the discussion of the curl}

  \begin{itemize}
  \tightlist
  \item
    The circulation of a vector field around a closed path
  \end{itemize}
\end{itemize}

\ELdisp{\text{Circulation of }\vect{A}\text{ around contour }C\triangleq\oint_C\vect{A}\cdot\dif\vect{\ell}}

\end{ELunit}

\begin{ELjoin}

\begin{itemize}
\tightlist
\item
  The physical meaning of the circulation depends on the kind of field that the vector \ELim{\vect{A}} represents.

  \begin{itemize}
  \tightlist
  \item
    If \ELim{\vect{A}} is the force acting on a body, the circulation of \ELim{\vect{A}} is the work done by the force in moving the body around the path.
  \end{itemize}
\end{itemize}

\begin{ELdefinition}{Curl}

The curl of a \textbf{vector field} \ELim{\vect{A}} at a point is a \textbf{vector} whose magnitude is the maximum net circulation of \ELim{\vect{A}} per unit area as the area about the point tends to zero, and whose direction is the direction of the normal to the area when the area is oriented so as to make the net circulation maximum.

\ELdisp{\operatorname{curl}\vect{A}\equiv\nabla\times\vect{A}\triangleq\lim_{\Delta s\to0}\frac{1}{\Delta s}\left[\uvec{n}\oint_C\vect{A}\cdot\dif\vect{\ell}\right]_{\max}}

\end{ELdefinition}

\end{ELjoin}

\begin{ELjoin}

\begin{itemize}
\tightlist
\item
  In Cartesian coordinates the curl is as follows
\end{itemize}

\begin{ELimportant}{}

\ELdisp{\nabla\times\vect{A}\ELbr =\begin{vmatrix}\uvec{x}&\uvec{y}&\uvec{z}\\[1mm]\dfrac{\partial}{\partial x}&\dfrac{\partial}{\partial y}&\dfrac{\partial}{\partial z}\\[3mm]A_x&A_y&A_z\end{vmatrix}}

\end{ELimportant}

\end{ELjoin}

\begin{ELjoin}

\begin{itemize}
\tightlist
\item
  The general expression of the curl
\end{itemize}

\begin{ELimportant}{}

\ELdisp{\nabla\times\vect{A}\ELbr =\frac{1}{h_1h_2h_3}\begin{vmatrix}\uvec{u_1}h_1&\uvec{u_2}h_2&\uvec{u_3}h_3\\[1mm]\dfrac{\partial}{\partial u_1}&\dfrac{\partial}{\partial u_2}&\dfrac{\partial}{\partial u_3}\\[3mm]h_1A_1&h_2A_2&h_3A_3\end{vmatrix}}

\end{ELimportant}

\end{ELjoin}

\begin{ELunit}

\begin{itemize}
\tightlist
\item
  \textbf{The concept of the curl}

  \begin{itemize}
  \tightlist
  \item
    The rotation of a paddle wheel in a stream of water with velocity vector \ELim{u}

    \begin{itemize}
    \tightlist
    \item
      The maximum angular velocity of the paddle wheel is proportional to the magnitude of the curl of the velocity vector.
    \item
      The axis of the wheel, by the right-hand rule, lies in the direction of the curl.
    \end{itemize}
  \end{itemize}
\end{itemize}

\ELfigure{m01-paddle-wheel}{}

\end{ELunit}

\ELnextnum{1-16}

\begin{ELexample}{}

Find the curl of the following vector
\ELdisp{\vect{A}\ELbr =\uvec{\phi}\left(\frac{k}{r}\right)}

\begin{ELsolution}

\ELdisp{\nabla\times\vect{A}\ELbr =\frac1r\begin{vmatrix}\uvec{r}&\uvec{\phi}r&\uvec{z}\\[1mm]\dfrac{\partial}{\partial r}&\dfrac{\partial}{\partial\phi}&\dfrac{\partial}{\partial z}\\[3mm]0&k&0\end{vmatrix}\ELbr =0}

\ELfigure{m01-ex16}{}

\end{ELsolution}

\end{ELexample}

\ELhold

\ELsection{Important Vector Theorems}{Important Vector Theorems}

\begin{ELunit}

\ELfigure{m01-vector-analysis-map-thm}{}

\end{ELunit}

\begin{ELtheorem}{Divergence theorem}

The volume integral of the divergence of a vector field equals the total outward flux of the vector through the surface enclosing the volume.
\ELdisp{\int_V\nabla\cdot\vect{A}\,dv\ELbr =\oint_S\vect{A}\cdot d\vect{s}}

\end{ELtheorem}

\begin{ELunit}

\ELfigure{m01-divergence-theorem}{}

\end{ELunit}

\begin{ELtheorem}{Stokes's theorem}

The surface integral of the curl of a vector field over an open surface equals the closed line integral of the vector along the contour bounding the surface.
\ELdisp{\int_S(\nabla\times\vect{A})\cdot d\vect{s}\ELbr =\oint_C\vect{A}\cdot\dif\vect{\ell}}

\end{ELtheorem}

\begin{ELunit}

\ELfigure{m01-stokes-theorem}{}

\end{ELunit}

\begin{ELtheorem}{Helmholtz's theorem}

A vector field is determined to within an additive constant if both its divergence and its curl are specified.

\end{ELtheorem}

\ELnextnum{1-17}

\begin{ELexample}{}

Verify the divergence theorem for the shell region enclosed by the spherical surfaces \ELim{R=R_1} and \ELim{R=R_2} \ELim{(R_2>R_1)} centred at the origin, for the following vector field.
\ELdisp{\vect{F}\ELbr =\uvec{R}kR}

\ELfigure{m01-ex17}{}

\begin{ELsolution}

\ELdisp{\int_V\nabla\cdot\vect{A}\,dv\ELbr =\oint_S\vect{A}\cdot d\vect{s}}

\begin{itemize}
\tightlist
\item
  For the outer surface
\end{itemize}

\ELdisp{R\ELbr =R_2,\qquad\ELbr  d\vect{s}\ELbr =\uvec{R}R_2^2\sin\theta\,d\theta\,d\phi}
\ELdisp{\int_{\substack{\text{outer}\\\text{surface}}}\vect{F}\cdot d\vect{s}\ELbr =\int_0^{2\pi}\!\!\int_0^\pi(kR_2)R_2^2\sin\theta\,d\theta\,d\phi\ELbr =4\pi kR_2^3}

\begin{itemize}
\tightlist
\item
  For the inner surface
\end{itemize}

\ELdisp{R\ELbr =R_1,\qquad\ELbr  d\vect{s}\ELbr =-\uvec{R}R_1^2\sin\theta\,d\theta\,d\phi}
\ELdisp{\int_{\substack{\text{inner}\\\text{surface}}}\vect{F}\cdot d\vect{s}\ELbr =-\int_0^{2\pi}\!\!\int_0^\pi(kR_1)R_1^2\sin\theta\,d\theta\,d\phi\ELbr =-4\pi kR_1^3}

\ELdisp{\oint_S\vect{F}\cdot d\vect{s}\ELbr =4\pi k(R_2^3-R_1^3)}
\ELdisp{\nabla\cdot\vect{F}\ELbr =\frac{1}{R^2}\frac{\partial}{\partial R}(R^2F_R)\ELbr =\frac{1}{R^2}\frac{\partial}{\partial R}(kR^3)\ELbr =3k}
\ELdisp{\int_V\nabla\cdot\vect{F}\,dv\ELbr =(\nabla\cdot\vect{F})V\ELbr =4\pi k(R_2^3-R_1^3)}

\end{ELsolution}

\end{ELexample}

\ELnextnum{1-18}

\begin{ELexample}{}

Verify Stokes's theorem over a quarter of a circular disk of radius 3 in the first quadrant for the following vector field.
\ELdisp{\vect{F}\ELbr =\uvec{x}xy-\uvec{y}2x}

\ELfigure{m01-ex18}{}

\begin{ELsolution}

\ELdisp{\int_S(\nabla\times\vect{A})\cdot d\vect{s}\ELbr =\oint_C\vect{A}\cdot\dif\vect{\ell}}
\ELdisp{\nabla\times\vect{F}\ELbr =\begin{vmatrix}\uvec{x}&\uvec{y}&\uvec{z}\\[1mm]\dfrac{\partial}{\partial x}&\dfrac{\partial}{\partial y}&\dfrac{\partial}{\partial z}\\[3mm]xy&-2x&0\end{vmatrix}\ELbr =-\uvec{z}(2+x)}
\ELdisp{\int_S(\nabla\times\vect{F})\cdot d\vect{s}\ELbr =\int_0^{\pi/2}\!\!\int_0^3-\uvec{z}(2+r\cos\phi)\cdot r\,dr\,d\phi\,\uvec{z}\ELbr =-9\left(1+\frac{\pi}{2}\right)}

\begin{itemize}
\tightlist
\item
  From \ELim{B} to \ELim{O}
\end{itemize}

\ELdisp{x\ELbr =0,\ \text{and}\ \vect{F}\cdot\dif\vect{\ell}\ELbr =\vect{F}\cdot(\uvec{y}\,dy)\ELbr =2x\,dy\ELbr =0}

\begin{itemize}
\tightlist
\item
  From \ELim{O} to \ELim{A}
\end{itemize}

\ELdisp{y\ELbr =0,\ \text{and}\ \vect{F}\cdot\dif\vect{\ell}\ELbr =\vect{F}\cdot(\uvec{x}\,dx)\ELbr =xy\,dx\ELbr =0}

\ELdisp{\oint_{ABOA}\vect{F}\cdot\dif\vect{\ell}\ELbr =\int_A^B\vect{F}\cdot\dif\vect{\ell}\ELbr =-9\left(1+\frac{\pi}{2}\right)}

(examined in Example \ELnumc{1-11})

\end{ELsolution}

\end{ELexample}

\ELhold

\ELsection{Important Vector Identities}{Important Vector Identities}

\begin{ELimportant}{Identity I}

The curl of the gradient of any scalar quantity is zero
\ELdisp{\nabla\times(\nabla V)\equiv0}

\end{ELimportant}

\begin{ELjoin}

\begin{itemize}
\tightlist
\item
  A curl-free vector field is called an \textbf{irrotational} or \textbf{conservative} field.
\item
  If a vector quantity is curl-free, it can be expressed as the gradient of a scalar quantity.
\end{itemize}

\begin{ELimportant}{Identity II}

The divergence of the curl of any vector field is zero.
\ELdisp{\nabla\cdot(\nabla\times\vect{A})\equiv0}

\end{ELimportant}

\end{ELjoin}

\begin{ELunit}

\begin{itemize}
\tightlist
\item
  A divergence-free vector field is called a \textbf{solenoidal} field.
\item
  If a vector quantity is divergence-free, it can be expressed as the curl of a vector quantity.
\end{itemize}

\end{ELunit}
